% ============================================================
%  04_kiem_thu_hop_den.tex  --  Phat hien loi bang AddressSanitizer
% ============================================================

\section{Phát hiện lỗi bằng AddressSanitizer}
\label{sec:hopden}

Kiểm thử hộp đen ở đây được thực hiện theo nghĩa: ta chạy chương trình với
đầu vào kích hoạt từng lỗi và quan sát \emph{hành vi đầu ra} của ASan mà không
cần đọc mã nguồn. Hai binary được biên dịch bằng lệnh:

\begin{lstlisting}[language=bash, caption={Biên dịch với AddressSanitizer}]
gcc -g -Wall -Wextra -fsanitize=address \
    -o c/vulnerable_inventory c/vulnerable_inventory.c
gcc -g -Wall -Wextra -fsanitize=address \
    -o c/secure_inventory    c/secure_inventory.c
\end{lstlisting}

% ────────────────────────────────────────────────────────────
\subsection{Kịch bản 1: Buffer Overflow (CWE-121)}

\subsubsection{Cách kích hoạt}

Chế độ \textit{overflow} truyền một chuỗi dài hơn 32 byte vào heap buffer:

\begin{lstlisting}[language=bash, caption={Kích hoạt heap buffer overflow}]
./c/vulnerable_inventory overflow "AAAAAAAAAAAAAAAAAAAAAAAAAAAAAAAAAAAAAAAA"
# payload: 39 ky tu 'A', vuot qua gioi han 32 byte
\end{lstlisting}

\subsubsection{Nguyên nhân trong mã nguồn}

Hàm \textit{add\_book\_title\_UNSAFE\_HEAP()} cấp phát đúng 32 byte
(\textit{TITLE\_MAX\_HEAP = 32}) nhưng gọi \textit{strcpy(buf, title)} mà không
kiểm tra \textit{strlen(title)}. Với payload 39 ký tự, 7 byte dư ghi vượt
ranh giới, đụng vào redzone của ASan.

\subsubsection{Thông báo lỗi ASan (rút gọn)}

\begin{lstlisting}[language=bash, caption={ASan output -- heap-buffer-overflow}]
=================================================================
==XXXXX==ERROR: AddressSanitizer: heap-buffer-overflow on address
    0x602000000030 at pc 0x... bp 0x... sp 0x...
WRITE of size 40 at 0x602000000030 thread T0
    #0 0x... in strcpy
    #1 0x... in add_book_title_UNSAFE_HEAP vulnerable_inventory.c:53
    #2 0x... in main vulnerable_inventory.c:81

0x602000000030 is located 0 bytes to the right of 32-byte region
    [0x602000000010, 0x602000000030)
allocated by thread T0 here:
    #0 0x... in malloc
    #1 0x... in add_book_title_UNSAFE_HEAP vulnerable_inventory.c:51

SUMMARY: AddressSanitizer: heap-buffer-overflow
         vulnerable_inventory.c:53 in add_book_title_UNSAFE_HEAP
\end{lstlisting}

\subsubsection{Phân tích}

ASan cho biết:
\begin{itemize}
    \item Loại lỗi: \textit{heap-buffer-overflow} -- ghi vượt ranh giới heap.
    \item Địa chỉ lỗi nằm ``0 bytes to the right of 32-byte region'' -- tức là
          ngay byte đầu tiên sau vùng 32 byte đã cấp phát.
    \item Stack trace: \textit{main} $\to$ \textit{add\_book\_title\_UNSAFE\_HEAP}
          $\to$ \textit{strcpy} (dòng 53).
    \item Tiến trình kết thúc với exit code khác 0.
\end{itemize}

\subsubsection{Kết quả bản secure}

\begin{lstlisting}[language=bash, caption={Bản secure -- overflow clean}]
./c/secure_inventory overflow "AAAAAAAAAAAAAAAAAAAAAAAAAAAAAAAAAAAAAAAA"
# Output:
Demo Heap Buffer Overflow (DA SUA), payload dai 39 ky tu:
  Tieu de (heap, an toan, co the bi cat bot): AAAAAAAAAAAAAAAAAAAAAAAAAAAAAAA
# Exit code: 0  |  Khong co ERROR trong stderr
\end{lstlisting}

\textit{snprintf} giới hạn ghi đúng 31 ký tự có nghĩa + 1 byte \textit{'\textbackslash0'},
chuỗi bị cắt an toàn, không có lỗi ASan.

% ────────────────────────────────────────────────────────────
\subsection{Kịch bản 2: Use-After-Free (CWE-416)}

\subsubsection{Cách kích hoạt}

\begin{lstlisting}[language=bash, caption={Kích hoạt use-after-free}]
./c/vulnerable_inventory uaf
\end{lstlisting}

\subsubsection{Nguyên nhân trong mã nguồn}

Hàm \textit{demo\_use\_after\_free()} gọi \textit{free(b)} ở dòng 66 rồi đọc
\textit{b->year} ở dòng 68. ASan đặt vùng nhớ đã giải phóng vào quarantine và
đánh dấu ``freed'' trong shadow memory, nên truy cập kế tiếp bị phát hiện ngay.

\subsubsection{Thông báo lỗi ASan (rút gọn)}

\begin{lstlisting}[language=bash, caption={ASan output -- heap-use-after-free}]
=================================================================
==XXXXX==ERROR: AddressSanitizer: heap-use-after-free on address
    0x602000000010 at pc 0x... bp 0x... sp 0x...
READ of size 4 at 0x602000000010 thread T0
    #0 0x... in demo_use_after_free vulnerable_inventory.c:68
    #1 0x... in main vulnerable_inventory.c:88

0x602000000010 is located 64 bytes inside of 72-byte region
    [0x602000000010, 0x602000000058)
freed by thread T0 here:
    #0 0x... in free
    #1 0x... in demo_use_after_free vulnerable_inventory.c:66

previously allocated by thread T0 here:
    #0 0x... in malloc
    #1 0x... in demo_use_after_free vulnerable_inventory.c:60

SUMMARY: AddressSanitizer: heap-use-after-free
         vulnerable_inventory.c:68 in demo_use_after_free
\end{lstlisting}

\subsubsection{Phân tích}

\begin{itemize}
    \item Loại lỗi: \textit{heap-use-after-free} -- đọc vùng nhớ đã giải phóng.
    \item ASan cho biết địa chỉ ``located 64 bytes inside of 72-byte region'' --
          trường \textit{year} nằm ở offset 64 trong struct (sau \textit{title}
          32 byte và \textit{author} 32 byte).
    \item Ba mốc thời gian được liệt kê: lúc cấp phát, lúc giải phóng, lúc truy
          cập lỗi -- cực kỳ hữu ích để debug.
    \item Stack trace: \textit{main} $\to$ \textit{demo\_use\_after\_free}
          (dòng 68).
\end{itemize}

\subsubsection{Kết quả bản secure}

\begin{lstlisting}[language=bash, caption={Bản secure -- UAF clean}]
./c/secure_inventory uaf
# Output:
Demo Use-After-Free (DA SUA):
  [FIXED] nam xuat ban (da sao chep truoc khi free): 2008
# Exit code: 0  |  Khong co ERROR trong stderr
\end{lstlisting}

Dữ liệu được sao chép vào biến \textit{year\_copy} trước khi \textit{free(b)},
con trỏ \textit{b} được đặt về \textit{NULL}. Không có lỗi ASan.

% ────────────────────────────────────────────────────────────
\subsection{Kịch bản 3: Memory Leak (CWE-401)}

\subsubsection{Cách kích hoạt}

\begin{lstlisting}[language=bash, caption={Kích hoạt memory leak với LeakSanitizer}]
ASAN_OPTIONS=detect_leaks=1 ./c/vulnerable_inventory leak
\end{lstlisting}

\subsubsection{Nguyên nhân trong mã nguồn}

Hàm \textit{demo\_memory\_leak(5)} chạy vòng lặp 5 lần, mỗi lần cấp phát
một đối tượng \textit{Book} (72 byte) gán cho con trỏ cục bộ \textit{b}.
Không có lời gọi \textit{free(b)} ở bất kỳ đâu trong hàm. Khi hàm kết thúc,
5 vùng nhớ (5 $\times$ 72 = 360 byte) bị mất hoàn toàn tham chiếu.

\subsubsection{Thông báo lỗi LeakSanitizer (rút gọn)}

\begin{lstlisting}[language=bash, caption={LeakSanitizer output -- memory leaks detected}]
Demo Memory Leak (CWE-401):
  (da cap phat 5 Book, KHONG giai phong -> memory leak)

=================================================================
==XXXXX==ERROR: LeakSanitizer: detected memory leaks

Direct leak of 360 byte(s) in 5 object(s) allocated from:
    #0 0x... in malloc
    #1 0x... in demo_memory_leak vulnerable_inventory.c:76
    #2 0x... in main vulnerable_inventory.c:93

SUMMARY: AddressSanitizer: 360 byte(s) leaked in 5 allocation(s).
\end{lstlisting}

\subsubsection{Phân tích}

\begin{itemize}
    \item Loại lỗi: \textit{detected memory leaks} -- tổng 360 byte rò rỉ.
    \item LeakSanitizer liệt kê chính xác số đối tượng (5 object), tổng byte (360)
          và vị trí cấp phát (dòng 76 trong \textit{demo\_memory\_leak}).
    \item Rò rỉ phân loại là ``Direct leak'' -- tham chiếu trực tiếp bị mất.
\end{itemize}

\subsubsection{Kết quả bản secure}

\begin{lstlisting}[language=bash, caption={Bản secure -- leak clean}]
ASAN_OPTIONS=detect_leaks=1 ./c/secure_inventory leak
# Output:
Demo Memory Leak (DA SUA):
  (da cap phat va giai phong day du 5 Book)
# Exit code: 0  |  Khong co "leaked" trong stderr
\end{lstlisting}

Tất cả 5 đối tượng và mảng con trỏ được giải phóng đầy đủ trong
\textit{demo\_memory\_leak\_FIXED()}. LeakSanitizer không báo lỗi.

% ────────────────────────────────────────────────────────────
\subsection{Bảng tổng hợp kết quả phát hiện lỗi}

\begin{table}[H]
\centering
\small
\begin{tabularx}{\textwidth}{|p{1.5cm}|p{1.8cm}|X|p{1.3cm}|}
\hline
\textbf{CWE} & \textbf{Kịch bản} & \textbf{Hàm gây lỗi / ASan report} & \textbf{Secure} \\ \hline
CWE-121 & \textit{overflow} &
  Hàm: \textit{add\_book\_title\_UNSAFE\_HEAP} \newline
  Lỗi: \textit{heap-buffer-overflow} -- ghi 40 byte vào vùng 32 byte & Sạch \\ \hline
CWE-416 & \textit{uaf} &
  Hàm: \textit{demo\_use\_after\_free} \newline
  Lỗi: \textit{heap-use-after-free} -- đọc sau \textit{free()} & Sạch \\ \hline
CWE-401 & \textit{leak} &
  Hàm: \textit{demo\_memory\_leak} \newline
  Lỗi: \textit{detected memory leaks} -- 360 byte / 5 objects & Sạch \\ \hline
\end{tabularx}
\caption{Tổng hợp 3 lỗi phát hiện bằng ASan trên bản vulnerable}
\end{table}

% ────────────────────────────────────────────────────────────
\subsection{Thí nghiệm bổ sung: tràn trong struct không bị ASan bắt}

\subsubsection{Mô tả thí nghiệm}

Đây là điểm sư phạm quan trọng nhất của bài thực hành. Hàm
\textit{add\_book\_title\_UNSAFE(\&b, payload)} dùng \textit{strcpy} ghi vào
\textit{b.title[32]} -- trường thuộc struct \textit{Book} được khai báo trên
stack. Nếu payload dài 35 ký tự, 3 byte dư ghi vào \textit{b.author[0..2]}.

\begin{lstlisting}[language=bash, caption={Tràn struct -- chế độ mặc định (safe\_input)}]
./c/vulnerable_inventory safe_input "AAAAAAAAAAAAAAAAAAAAAAAAAAAAAAAAAAA"
# payload: 35 ky tu -> tran 3 byte vao b.author
# Ket qua: KHONG co bao loi ASan, exit code 0
\end{lstlisting}

\subsubsection{Giải thích kỹ thuật}

Bố trí bộ nhớ của \textit{Book b} trên stack:

\begin{lstlisting}[language=bash, caption={Bố trí bộ nhớ struct Book trên stack}]
[ redzone | title[0..31] | author[0..31] | year | copies | redzone ]
              ^                                              ^
              |__ ASan chi dat redzone o day _______________|
                      (bien ngoai cung cua allocation)
\end{lstlisting}

Ghi 35 byte vào \textit{title}: byte 0--31 vào \textit{title}, byte 32--34 vào
\textit{author[0..2]}. Tất cả vẫn nằm trong allocation của struct -- không
đụng đến redzone -- nên ASan không phát hiện. Tuy nhiên, \textit{b.author}
đã bị ghi đè bởi dữ liệu tấn công.

\subsubsection{So sánh với heap buffer riêng}

Ngược lại, \textit{add\_book\_title\_UNSAFE\_HEAP()} cấp phát \textit{malloc(32)}
riêng biệt. ASan đặt redzone \emph{ngay sau} byte thứ 31. Mọi ghi từ byte thứ
32 trở đi đụng ngay redzone $\Rightarrow$ ASan phát hiện tức thì.

\begin{table}[H]
\centering
\begin{tabularx}{\textwidth}{|X|p{2.5cm}|p{4cm}|}
\hline
\textbf{Tình huống} & \textbf{ASan phát hiện?} & \textbf{Lý do} \\ \hline
Tràn trong struct trên stack & KHÔNG & Ghi vào field kề bên, vẫn trong cùng allocation; không đụng redzone \\ \hline
Tràn heap buffer độc lập & CÓ & Redzone nằm ngay sau biên vùng \textit{malloc(32)}; bị đụng ngay lập tức \\ \hline
\end{tabularx}
\caption{So sánh khả năng phát hiện của ASan theo bố trí bộ nhớ}
\end{table}

\subsubsection{Kết luận thực nghiệm}

Kết quả thí nghiệm này minh chứng trực tiếp rằng ASan không phải công
cụ chứng minh chương trình an toàn. Nó chỉ phát hiện lỗi trong các kịch bản
được kiểm tra, với bố trí bộ nhớ cụ thể. Lỗi tràn bộ đệm thực sự có thể tồn
tại và gây hại mà không bao giờ sinh ra thông báo ASan -- đặc biệt khi tràn
xảy ra giữa các trường bên trong cùng một struct.

Đây là lý do tại sao cần kết hợp nhiều lớp bảo vệ: kiểm tra độ dài đầu vào,
dùng API an toàn (\textit{snprintf}), code review, phân tích tĩnh, và kiểm
thử với nhiều kịch bản edge-case.
